% GENERATED FILE. Edit canonical lessons, then run npm run generate:books.
% canonical-source-hash: fnv1a64:a4a88d2f115fe542
% canonical-lessons: RU-C333-solit, RU-C333-perchit, RU-C333-posypat, RU-C333-priprava, RU-C333-marinovat

\chapter{Seasoning}
\label{ch:ru-seasoning}
\hlchaptermodality{333}

\begin{chapteropening}
\textbf{By the end of this chapter:} \emph{I can describe seasoning food: salting, peppering, sprinkling, a seasoning and marinating.}

Complete the last lesson of chapter 333: I can describe seasoning food: salting, peppering, sprinkling, a seasoning and marinating.
\end{chapteropening}

\section[\texorpdfstring{\ru{солить} (solít)}{солить (solít)}]{\ru{солить} (solít) — to salt}

\label{lesson:RU-C333-solit}



\begin{quote}
Before the new one: say the Russian for to get burnt, then the Russian for to brown.
\end{quote}

\subsection*{You'll want to know: \ru{солить}}
\textbf{\ru{солить}} (\emph{solít}) — \textquotedblleft{}to salt\textquotedblright{}.

\begin{quote}
\ru{Ты} \ru{уже} \ru{солил} \ru{суп}?
\end{quote}

\begin{quote}
\emph{Ty uzhé solíl sup?}
\end{quote}

\begin{quote}
Have you salted the soup yet?
\end{quote}

Imperfective; perfective \emph{\ru{посолить}} (\emph{posolít}). Also to pickle in salt: \emph{\ru{солить} \ru{огурцы}}.

\begin{cousinweb}
\emph{\ru{солить}} is built on \emph{\ru{соль}} \textquotedblleft{}salt\textquotedblright{}, an inherited Slavic word from the PIE root of Latin \emph{sal} and English \emph{salt}; EtymDB lists no parent for the verb.
\end{cousinweb}

\begin{grammarlens}[title={What you've built}]
One more word for telling, explaining and describing at B1.
\end{grammarlens}

\subsection*{Your turn}
\begin{itemize}
\raggedright
  \item \emph{Say it:} \emph{solít}
  \item \emph{Say it:} \emph{solít}, once more
  \item \emph{Say it:} a sentence of your own with \emph{solít}
  \item \emph{Recall:} say the Russian for to get burnt, then the Russian for to brown, then say \emph{solít} again
\end{itemize}

\begin{tcolorbox}[breakable,colback=teal!4,colframe=teal!35!black,title={Before you move on}]
What does \ru{солить} mean? (\textquotedblleft{}To salt\textquotedblright{}.) Say it once more.
\end{tcolorbox}
\section[\texorpdfstring{\ru{перчить} (perchít)}{перчить (perchít)}]{\ru{перчить} (perchít) — to pepper}

\label{lesson:RU-C333-perchit}



\begin{quote}
Before the new one: say the Russian for to brown, then the Russian for to salt.
\end{quote}

\subsection*{You'll want to know: \ru{перчить}}
\textbf{\ru{перчить}} (\emph{perchít}) — \textquotedblleft{}to pepper\textquotedblright{}.

\begin{quote}
\ru{Он} \ru{всегда} \ru{много} \ru{перчит} \ru{мясо}.
\end{quote}

\begin{quote}
\emph{On vsegdá mnógo perchít myáso.}
\end{quote}

\begin{quote}
He always puts a lot of pepper on meat.
\end{quote}

Imperfective; perfective \emph{\ru{поперчить}} (\emph{poperchít}).

\begin{cousinweb}
\emph{\ru{перчить}} is built on \emph{\ru{перец}} \textquotedblleft{}pepper\textquotedblright{}, inherited through Old East Slavic \emph{pĭpĭrĭcĭ} from Proto-Slavic \emph{pĭpĭrĭ}, which EtymDB traces to Latin \emph{piper} \textquotedblleft{}pepper\textquotedblright{}.
\end{cousinweb}

\begin{grammarlens}[title={What you've built}]
One more word for telling, explaining and describing at B1.
\end{grammarlens}

\subsection*{Your turn}
\begin{itemize}
\raggedright
  \item \emph{Say it:} \emph{perchít}
  \item \emph{Say it:} \emph{perchít}, once more
  \item \emph{Say it:} a sentence of your own with \emph{perchít}
  \item \emph{Recall:} say the Russian for to brown, then the Russian for to salt, then say \emph{perchít} again
\end{itemize}

\begin{tcolorbox}[breakable,colback=teal!4,colframe=teal!35!black,title={Before you move on}]
What does \ru{перчить} mean? (\textquotedblleft{}To pepper\textquotedblright{}.) Say it once more.
\end{tcolorbox}
\section[\texorpdfstring{\ru{посыпать} (posýpat)}{посыпать (posýpat)}]{\ru{посыпать} (posýpat) — to sprinkle}

\label{lesson:RU-C333-posypat}



\begin{quote}
Before the new one: say the Russian for to salt, then the Russian for to pepper.
\end{quote}

\subsection*{You'll want to know: \ru{посыпать}}
\textbf{\ru{посыпать}} (\emph{posýpat}) — \textquotedblleft{}to sprinkle\textquotedblright{}.

\begin{quote}
\ru{Она} \ru{посыпала} \ru{салат} \ru{укропом}.
\end{quote}

\begin{quote}
\emph{Oná posýpala salát ukrópom.}
\end{quote}

\begin{quote}
She sprinkled the salad with dill.
\end{quote}

Perfective, stressed on \emph{-\ru{сы}-}: \emph{\ru{посыплю}}; the imperfective is \emph{\ru{посыпать}} (\emph{posypát}). The topping is in the instrumental.

\begin{cousinweb}
\emph{\ru{посыпать}} is built on \emph{\ru{по}-} and \emph{\ru{сыпать}} \textquotedblleft{}to pour grains or powder\textquotedblright{}, an inherited Slavic verb; EtymDB lists no parent.
\end{cousinweb}

\begin{grammarlens}[title={What you've built}]
One more word for telling, explaining and describing at B1.
\end{grammarlens}

\subsection*{Your turn}
\begin{itemize}
\raggedright
  \item \emph{Say it:} \emph{posýpat}
  \item \emph{Say it:} \emph{posýpat}, once more
  \item \emph{Say it:} a sentence of your own with \emph{posýpat}
  \item \emph{Recall:} say the Russian for to salt, then the Russian for to pepper, then say \emph{posýpat} again
\end{itemize}

\begin{tcolorbox}[breakable,colback=teal!4,colframe=teal!35!black,title={Before you move on}]
What does \ru{посыпать} mean? (\textquotedblleft{}To sprinkle\textquotedblright{}.) Say it once more.
\end{tcolorbox}
\section[\texorpdfstring{\ru{приправа} (pripráva)}{приправа (pripráva)}]{\ru{приправа} (pripráva) — a seasoning}

\label{lesson:RU-C333-priprava}



\begin{quote}
Before the new one: say the Russian for to pepper, then the Russian for to sprinkle.
\end{quote}

\subsection*{You'll want to know: \ru{приправа}}
\textbf{\ru{приправа}} (\emph{pripráva}) — \textquotedblleft{}a seasoning\textquotedblright{}.

\begin{quote}
\ru{Эта} \ru{приправа} \ru{очень} \ru{острая}.
\end{quote}

\begin{quote}
\emph{Éta pripráva óchen óstraya.}
\end{quote}

\begin{quote}
This seasoning is very hot.
\end{quote}

Feminine.

\begin{cousinweb}
\emph{\ru{приправа}} is built on \emph{\ru{приправить}} \textquotedblleft{}to season\textquotedblright{}, from \emph{\ru{при}-} and \emph{\ru{править}} \textquotedblleft{}to set right\textquotedblright{}, from \emph{\ru{правый}} \textquotedblleft{}right\textquotedblright{}: what puts a dish right. EtymDB lists no parent.
\end{cousinweb}

\begin{grammarlens}[title={What you've built}]
One more word for telling, explaining and describing at B1.
\end{grammarlens}

\subsection*{Your turn}
\begin{itemize}
\raggedright
  \item \emph{Say it:} \emph{pripráva}
  \item \emph{Say it:} \emph{pripráva}, once more
  \item \emph{Say it:} a sentence of your own with \emph{pripráva}
  \item \emph{Recall:} say the Russian for to pepper, then the Russian for to sprinkle, then say \emph{pripráva} again
\end{itemize}

\begin{tcolorbox}[breakable,colback=teal!4,colframe=teal!35!black,title={Before you move on}]
What does \ru{приправа} mean? (\textquotedblleft{}A seasoning\textquotedblright{}.) Say it once more.
\end{tcolorbox}
\section[\texorpdfstring{\ru{мариновать} (marinovát)}{мариновать (marinovát)}]{\ru{мариновать} (marinovát) — to marinate; to pickle}

\label{lesson:RU-C333-marinovat}



\begin{quote}
Before the new one: say the Russian for to sprinkle, then the Russian for a seasoning.
\end{quote}

\subsection*{You'll want to know: \ru{мариновать}}
\textbf{\ru{мариновать}} (\emph{marinovát}) — \textquotedblleft{}to marinate; to pickle\textquotedblright{}.

\begin{quote}
\ru{Мясо} \ru{для} \ru{шашлыка} \ru{маринуют} \ru{всю} \ru{ночь}.
\end{quote}

\begin{quote}
\emph{Myáso dlya shashlyká marinúyut vsyu noch.}
\end{quote}

\begin{quote}
The meat for kebabs is marinated all night.
\end{quote}

Imperfective; the perfective is \emph{\ru{замариновать}} (\emph{zamarinovát}). Present \emph{\ru{мариную}}.

\begin{cousinweb}
\emph{\ru{мариновать}} is usually traced through French \emph{mariner} \textquotedblleft{}to pickle in brine\textquotedblright{}, built on \emph{marin} \textquotedblleft{}of the sea\textquotedblright{}, from Latin \emph{mare} \textquotedblleft{}sea\textquotedblright{}; EtymDB lists no parent.
\end{cousinweb}

\begin{grammarlens}[title={What you've built}]
One more word for telling, explaining and describing at B1.
\end{grammarlens}

\subsection*{Your turn}
\begin{itemize}
\raggedright
  \item \emph{Say it:} \emph{marinovát}
  \item \emph{Say it:} \emph{marinovát}, once more
  \item \emph{Say it:} a sentence of your own with \emph{marinovát}
  \item \emph{Recall:} say the Russian for to sprinkle, then the Russian for a seasoning, then say \emph{marinovát} again
\end{itemize}

\begin{tcolorbox}[breakable,colback=teal!4,colframe=teal!35!black,title={Before you move on}]
What does \ru{мариновать} mean? (\textquotedblleft{}To marinate; to pickle\textquotedblright{}.) Say it once more.
\end{tcolorbox}
